%% toc: P3 33 33.5,33.6,33.7,33.8
\section{Framing, Saturation and Footnoted Results in Morley's \emph{Categoricity in Power}}
\label{sec:p03-w03-p03-034-morley-framing}

This section supplies the parts of Morley's 1965 article that surround the
main line of \cref{sec:p03-n-a-006-01}: the opening problem and its two
classical examples, the count of the canonical basis of a type space, the
sense in which the global derivative generalizes the Cantor--Bendixson
derivative, and the footnoted results on saturation that Morley reports
from Vaught, Keisler and the Morley--Vaught paper
\cite[pp.~514--516, 518, 520, 522, 535]{Morley1965}. The inputs are
Morley's standing data of \cref{def:p03-morley-standing}: a complete
theory \(\Sigma\) with an infinite model in an at most countable language,
its definitional relational expansion \(T=\Sigma^*\), the class
\(\mathcal N(T)\) of induced subsystems of models of \(T\), and
monomorphisms, which preserve and reflect every relation and are
injective. The conclusions are these. Dense linear order is categorical in
\(\aleph_0\) and in no uncountable power
(\cref{thm:p03-w03-p03-034-dlo-powers}). The canonical basis of
\(S(A)\) has power \(\kappa(A)\) modulo \(\aleph_0\)
(\cref{thm:p03-w03-p03-034-basis-power}). Morley's derivative and its
closedness and image theorem hold in every category of continuous onto
maps between compact Hausdorff spaces with identities and amalgamation,
and the one-space case is the Cantor--Bendixson derivative
(\cref{thm:p03-w03-p03-034-abstract-derivative,thm:p03-w03-p03-034-cantor-bendixson}).
A system in \(\mathcal N(T)\) is saturated exactly when it is universal and
homogeneous in Jónsson's sense
(\cref{thm:p03-w03-p03-034-universal-homogeneous}). Categoricity in an
increasing sequence of powers gives categoricity in their supremum
without any continuum hypothesis
(\cref{thm:p03-w03-p03-034-limit-power}). Models saturated over countable
subsystems exist in every power \(\kappa=\kappa^{\aleph_0}\) without total
transcendence, and a theory that is not totally transcendental has none of
power below \(2^{\aleph_0}\)
(\cref{thm:p03-w03-p03-034-mv-existence,thm:p03-w03-p03-034-continuum-bound}).
None of these results uses the later strongly minimal analysis of
\cref{sec:p03-n-a-006-02}.

Throughout, \(\kappa(X)\) is the cardinality of a set \(X\) and
\(\kappa(A)=\kappa(|A|)\) is the power of a system \(A\); a bare \(\kappa\)
or \(\lambda\) is a cardinal. As in the source, \(\kappa=\lambda\)
\emph{modulo} \(\aleph_0\) means \(\kappa+\aleph_0=\lambda+\aleph_0\)
\cite[p.~515]{Morley1965}. For \(A\subseteq B\) in \(\mathcal N(T)\)
with \(B\) contained in a model, an element \(b\) of that model
\emph{realizes} the point of \(S(A)\) consisting of the formulas of
\(F(A)\) that \(b\) satisfies. For an isomorphism \(f:D\to E\) between
systems, \(f^*:S(E)\to S(D)\) is a bijection, with inverse
\((f^{-1})^*\), by functoriality of substitution; for \(p\in S(D)\) the
point \((f^{-1})^*(p)\in S(E)\) is called the \emph{transport} of \(p\)
along \(f\).

\subsection{The opening problem and its two examples}
\label{sec:p03-w03-p03-034-opening}

Morley opens with the history of the notion. Categoricity in power was
introduced by Łoś and by Vaught in 1954 (his references [9] and [16]), and
they observed at once that the two directions of transfer between
\(\aleph_0\) and the uncountable powers both fail: dense linear order
without endpoints is categorical in \(\aleph_0\) but in no higher power,
while the theory of algebraically closed fields of characteristic \(0\) is
categorical in every uncountable power but not in \(\aleph_0\)
\cite[p.~514]{Morley1965}. Łoś then asked whether a theory categorical in
one uncountable power must be categorical in every uncountable power. The
article answers this affirmatively and proves more: under categoricity in
some uncountable power, every uncountable model is saturated
(\cref{thm:p03-morley-original-categoricity}). Morley glosses
\emph{saturated} informally as \emph{universal-homogeneous}: a model
containing an element of every possible elementary type over each of its
subsystems of smaller power \cite[p.~514]{Morley1965}. The precise
equivalence of the two notions in \(\mathcal N(T)\) is
\cref{thm:p03-w03-p03-034-universal-homogeneous} below.

A theory has, up to isomorphism, at most one saturated model in each power;
this is Morley's Theorem~5.1 (\cref{thm:p03-morley-saturation-existence}).
In the introduction he records in a footnote that the Morley--Vaught paper
(his reference [10]) had derived it from a more general theorem of Jónsson
(his reference [5]) \cite[p.~514, n.~2]{Morley1965}. A further footnote
states that, apart from minor emendations, the article is Morley's
doctoral dissertation submitted to the University of Chicago in August
1962 \cite[p.~514, n.~1]{Morley1965}. The completeness of a theory that is
categorical in an infinite power and has no finite models, used to
normalize the standing theory, is attributed there to Vaught and Łoś
\cite[p.~516]{Morley1965}.

The first half of the opening history is the following theorem; the second
half is \cref{thm:p03-morley-field-example}, which shows that every
uncountable algebraically closed field of fixed characteristic is
determined by its power, while the countable ones have the distinct
transcendence degrees \(0,1,2,\ldots,\aleph_0\).

\begin{hbtheorem}{thm:p03-w03-p03-034-dlo-powers}{proved=annex:p03-w03-p03-034-dlo-powers}{Dense linear order is categorical in \(\aleph_0\) and in no higher power}
Let \(\mathsf{DLO}\) be the theory of dense linear orders without
endpoints. Any two countable models of \(\mathsf{DLO}\) are isomorphic.
For every uncountable cardinal \(\kappa\), let \(L_\kappa\) be the set
\(\kappa\times\mathbb Q\) ordered lexicographically, with the first
coordinate dominant, and let \(L_\kappa^{\mathrm{rev}}\) be the same set
with the reverse order. Both are models of \(\mathsf{DLO}\) of power
\(\kappa\); \(L_\kappa\) contains a strictly increasing sequence of
order type \(\omega_1\), and \(L_\kappa^{\mathrm{rev}}\) does not.
Hence \(\mathsf{DLO}\) is not categorical in any uncountable power.
\end{hbtheorem}

The countable case is Cantor's back-and-forth theorem in the form of
\cref{thm:dlo-countable-categorical}. Morley returns to dense linear order
after his Theorem~3.8: it is categorical in \(\aleph_0\) but not totally
transcendental, because the Dedekind cuts of the rational order give
continuum many points of \(S(A)\) for a countable \(A\)
\cite[p.~528]{Morley1965}; this is the final clause of
\cref{thm:p03-morley-thin-models}.

\subsection{The canonical basis of a type space}
\label{sec:p03-w03-p03-034-basis}

After defining the Stone space \(S(A)\) of the formula algebra \(F(A)\),
Morley notes that \(S(A)\) has a basis of power \(\kappa(A)\) modulo
\(\aleph_0\) \cite[p.~518]{Morley1965}. The basis meant is the canonical
one, \(\{U_\varphi:\varphi\in F(A)\}\) of
\cref{def:p03-morley-type-space}. \Cref{thm:p03-morley-types} gives the
upper bound; the lower bound, which makes the count an equality modulo
\(\aleph_0\), is the following.

\begin{hblemma}{thm:p03-w03-p03-034-basis-power}{proved=annex:p03-w03-p03-034-basis-power}{The canonical basis has power \(\kappa(A)\) modulo \(\aleph_0\)}
For every \(A\in\mathcal N(T)\), the map \(\varphi\mapsto U_\varphi\) is a
bijection from \(F(A)\) onto the canonical clopen basis of \(S(A)\), and
\[
 \kappa(A)\leq\kappa(F(A))\leq\kappa(A)+\aleph_0 ,
 \qquad
 \kappa(\{U_\varphi:\varphi\in F(A)\})+\aleph_0=\kappa(A)+\aleph_0 .
\]
The modulus cannot be dropped: if \(\Sigma\) is the theory of an infinite
set with equality only and \(A=\varnothing\), the canonical basis has
exactly two members.
\end{hblemma}

\subsection{The global derivative as an analogue of the Cantor--Bendixson derivative}
\label{sec:p03-w03-p03-034-analogue}

Morley states that his results depend essentially on an analogue of the
analysis of a topological space by its derived spaces and the
Cantor--Bendixson theorem \cite[p.~514]{Morley1965}. The derivative
\(S^\alpha(A)\) of \cref{def:p03-morley-global-rank} is analogous to,
but different from, the sequence of derived spaces of one space: a point is
removed at stage \(\alpha\) only if every point over it, in every space of
the category, is isolated at that stage \cite[pp.~514--515]{Morley1965}.
He calls Theorem~2.8 (\cref{thm:p03-morley-countability}: total
transcendence is equivalent to countability of \(S(A)\) for all countable
\(A\)) an analogue of the Cantor--Bendixson theorem, proved in a similar
way \cite[pp.~515 and 522]{Morley1965}. In the classical theorem a closed
subset of a Polish space is countable exactly when its perfect kernel is
empty, and a nonempty perfect kernel is split repeatedly into two disjoint
nonempty closed pieces to produce continuum many points. In Morley's
proof a nonempty residual space \(S^{\alpha_T}\) is split in the same
binary way, over a growing chain of finite systems, to produce \(2^{\aleph_0}\)
points over a countable system.

Morley also remarks, without proof, that it will be evident that Definition~2.2
and many of the subsequent results and proofs remain valid in many other
categories of continuous onto maps between compact spaces
\cite[p.~520]{Morley1965}. He does not list the categories or the
results. The following definition and theorem make the remark precise for
Definition~2.2 and Theorem~2.3, using exactly the properties that his proof
of Theorem~2.3 invokes: compactness of the spaces, the Hausdorff property of
the codomains, continuity and surjectivity of the maps, identities and
composition, and the dual of the amalgamation property~1.2(d)
\cite[pp.~520--521]{Morley1965}. Which of the later results of the
article survive in which categories is not specified in the source and is
not decided here.

\begin{hbdefinition}{def:p03-w03-p03-034-abstract-category}{A derivative category and its derivative}
A \emph{derivative category} \(\mathcal C\) is a class of compact
Hausdorff spaces together with a class of continuous maps of these spaces
onto one another, written \(\phi:Y\to X\), such that every identity map
belongs to \(\mathcal C\), the composite of two maps of \(\mathcal C\)
belongs to \(\mathcal C\), and \(\mathcal C\) has \emph{amalgamation}: for
maps \(\phi:Y\to X\) and \(\psi:V\to X\) of \(\mathcal C\) there are a
space \(W\) and maps \(\phi':W\to Y\) and \(\psi':W\to V\) of
\(\mathcal C\) with \(\phi\circ\phi'=\psi\circ\psi'\). Simultaneously for
all spaces \(X\) of \(\mathcal C\) and by recursion on the ordinal
\(\alpha\), put
\[
 X^\alpha=X\setminus\bigcup_{\beta<\alpha}\operatorname{Tr}^\beta(X),
\]
and let \(p\in\operatorname{Tr}^\alpha(X)\) when \(p\in X^\alpha\) and,
for every map \(\phi:Y\to X\) of \(\mathcal C\), every point of
\(\phi^{-1}(p)\cap Y^\alpha\) is isolated in \(Y^\alpha\).
\end{hbdefinition}

For Morley's category, the spaces are the \(S(A)\) with
\(A\in\mathcal N(T)\) and the maps are the restriction maps \(f^*\) of
monomorphisms \(f\); identities and composites are again of this form, the
spaces are compact Hausdorff and the maps continuous onto by
\cref{thm:p03-morley-types}, and amalgamation is the dual of the
amalgamation of \(\mathcal N(T)\) in \cref{thm:p03-morley-normalization}.
In that case \(X^\alpha\) and \(\operatorname{Tr}^\alpha(X)\) are the
\(S^\alpha(A)\) and \(\operatorname{Tr}^\alpha(A)\) of
\cref{def:p03-morley-global-rank}.

\begin{hbtheorem}{thm:p03-w03-p03-034-abstract-derivative}{proved=annex:p03-w03-p03-034-abstract-derivative}{Morley's Theorem~2.3 in a derivative category}
Let \(\mathcal C\) be a derivative category, \(X\) a space of
\(\mathcal C\) and \(\alpha\) an ordinal.
\begin{enumerate}
\item \(X^\alpha\) is closed in \(X\), hence compact, and
\(X^{\alpha+1}=X^\alpha\setminus\operatorname{Tr}^\alpha(X)\).
\item For every map \(\phi:Y\to X\) of \(\mathcal C\),
\(\phi(Y^\alpha)=X^\alpha\).
\item For every such map and every \(p\in X^\alpha\),
\(p\in\operatorname{Tr}^\alpha(X)\) if and only if
\(\phi^{-1}(p)\cap Y^\alpha\subseteq\operatorname{Tr}^\alpha(Y)\).
\end{enumerate}
\end{hbtheorem}

\begin{hbcorollary}{thm:p03-w03-p03-034-cantor-bendixson}{proved=annex:p03-w03-p03-034-cantor-bendixson}{The one-space category gives the Cantor--Bendixson derivative}
Let \(X\) be a compact Hausdorff space and let \(\mathcal C\) have \(X\)
as its only space and the identity of \(X\) as its only map. Then
\(\mathcal C\) is a derivative category, \(\operatorname{Tr}^\alpha(X)\)
is the set of isolated points of \(X^\alpha\), and \(X^\alpha\) is the
\(\alpha\)-th Cantor--Bendixson derivative \(X^{(\alpha)}\) of \(X\),
defined by \(X^{(0)}=X\), \(X^{(\alpha+1)}\) the set of non-isolated
points of \(X^{(\alpha)}\), and intersections at limits.
\end{hbcorollary}

The corollary is the precise form of the analogy in the one-space case.
Morley's Example~III shows the converse embedding: there the residual space
\(S^1(A)\) is homeomorphic to a closed set \(Y\subseteq2^\omega\) and its
global derivatives are the ordinary derived sets of \(Y\)
(\cref{thm:p03-morley-unary-example}). Neither result says that
\(S^\alpha(A)\) is the Cantor--Bendixson derivative of \(S(A)\) in general:
in Morley's category the fibre condition is tested over every extension of
\(A\), and an isolated point of \(S(A)\) need not be algebraic
(\cref{def:p03-morley-global-rank}).

\subsection{Saturation and Jónsson's universal homogeneous systems}
\label{sec:p03-w03-p03-034-universal-homogeneous}

Morley's Theorem~5.1, uniqueness of saturated models of a given power, is
taken from the Morley--Vaught paper, which works with universal homogeneous
systems, a notion and terminology of Jónsson. In a footnote Morley gives the
definitions, reports that Jónsson proved under certain simple conditions on
a class that any two universal homogeneous members of the same power are
isomorphic, and states that for the class \(\mathcal N(T)\) universal
homogeneous is equivalent to saturated, an equivalence shown by Vaught in
the countable case (his reference [18]) and by Keisler in the uncountable
case (Theorem~A2 of his reference [8]) \cite[p.~535, n.~22]{Morley1965}.
The countable case is Vaught's Theorem~4.6 in
\cref{thm:p03-vaught-saturation-characterization}. The theorem below proves
the equivalence for every infinite power directly in Morley's setting; it
does not reconstruct Keisler's argument.

\begin{hbdefinition}{def:p03-w03-p03-034-universal-homogeneous}{Universal and homogeneous systems}
Let \(K\) be a class of relational systems of one similarity type and
\(A\in K\). Following Morley's footnote \cite[p.~535, n.~22]{Morley1965}:
\(A\) is \emph{universal for} \(K\) when every \(B\in K\) with
\(\kappa(B)\leq\kappa(A)\) has a monomorphism into \(A\), that is, \(A\)
contains an isomorphic image of \(B\); and \(A\) is \emph{homogeneous in}
\(K\) when, for all \(B_1,B_2\in K\) with \(B_1,B_2\subseteq A\) and
\(\kappa(B_i)<\kappa(A)\), every isomorphism \(f:B_1\to B_2\) extends to
an automorphism of \(A\). Here \(K=\mathcal N(T)\), so every subsystem of
\(A\) belongs to \(K\).
\end{hbdefinition}

\begin{hbtheorem}{thm:p03-w03-p03-034-universal-homogeneous}{proved=annex:p03-w03-p03-034-universal-homogeneous}{Saturated is universal homogeneous in \(\mathcal N(T)\)}
Let \(A\in\mathcal N(T)\) be infinite. Then \(A\) is saturated
(\cref{def:p03-morley-saturation}) if and only if \(A\) is universal for
\(\mathcal N(T)\) and homogeneous in \(\mathcal N(T)\).
\end{hbtheorem}

Combined with \cref{thm:p03-morley-saturation-existence}, the theorem
shows that Morley's Theorem~5.1 and Jónsson's uniqueness of universal
homogeneous systems are, for \(\mathcal N(T)\), the same statement.

Morley adds in a further footnote that Vaught had noticed that categoricity
in power could be approached through saturated models, since a theory all
of whose models of power \(\kappa\) are saturated is categorical in
\(\kappa\); Vaught proved, assuming the generalized continuum hypothesis,
that a theory categorical in an increasing sequence of powers is
categorical in the limit power (Morley's references [10] and [17])
\cite[p.~535, n.~23]{Morley1965}. For Morley's standing theories the
conclusion follows from the main theorem with no hypothesis on cardinal
arithmetic.

\begin{hbcorollary}{thm:p03-w03-p03-034-limit-power}{proved=annex:p03-w03-p03-034-limit-power}{Vaught's limit-power statement without the continuum hypothesis}
Let \(\Sigma\) be a complete theory with an infinite model in an at most
countable language. Let \((\kappa_\xi)_{\xi<\gamma}\) be a strictly
increasing sequence of infinite cardinals indexed by a nonzero limit ordinal
\(\gamma\), with supremum \(\lambda\). If \(\Sigma\) is categorical in
every \(\kappa_\xi\), then \(\Sigma\) is categorical in \(\lambda\).
\end{hbcorollary}

\subsection{Countably saturated models without total transcendence}
\label{sec:p03-w03-p03-034-countable-saturation}

Morley's Theorem~5.2 gives, for a totally transcendental theory and every
\(\kappa>\aleph_0\), a model of power \(\kappa\) saturated over countable
subsystems (\cref{thm:p03-morley-saturation-existence}). His footnote
records two further facts \cite[p.~535, n.~24]{Morley1965}. First, in the
case \(\kappa=\kappa^{\aleph_0}\) the conclusion had been proved in the
Morley--Vaught paper without total transcendence. Second, he asserts that
one can give an example of a theory that is not totally transcendental and a
cardinal \(\kappa>\aleph_0\) with \(\kappa^{\aleph_0}\neq\kappa\) such
that no model of the theory of power \(\kappa\) is countably saturated; the
article does not give the example. Here \emph{countably saturated} is read
as saturated over countable subsystems, the notion of
\cref{def:p03-morley-saturation}.

\begin{hbtheorem}{thm:p03-w03-p03-034-mv-existence}{proved=annex:p03-w03-p03-034-mv-existence}{Saturation over countable subsystems when \(\kappa^{\aleph_0}=\kappa\)}
Let \(\kappa>\aleph_0\) with \(\kappa^{\aleph_0}=\kappa\). Then \(T\) has
a model of power \(\kappa\) that is saturated over countable subsystems.
Total transcendence is not assumed.
\end{hbtheorem}

\begin{hbtheorem}{thm:p03-w03-p03-034-continuum-bound}{proved=annex:p03-w03-p03-034-continuum-bound}{A lower bound for countably saturated models}
If \(T\) is not totally transcendental, every model of \(T\) that is
saturated over countable subsystems has power at least \(2^{\aleph_0}\).
Consequently, if \(2^{\aleph_0}>\aleph_1\), the cardinal
\(\kappa=\aleph_1\) satisfies \(\kappa^{\aleph_0}\neq\kappa\), and no
theory that is not totally transcendental, for instance \(\mathsf{DLO}\),
has a model of power \(\kappa\) saturated over countable subsystems.
\end{hbtheorem}

\begin{hbwarning}{def:p03-w03-p03-034-example-limit}{What is and is not supplied}
\Cref{thm:p03-w03-p03-034-continuum-bound} exhibits the situation Morley
asserts only under the additional hypothesis \(2^{\aleph_0}>\aleph_1\).
Under the continuum hypothesis there is no uncountable
\(\kappa<2^{\aleph_0}\), so an example must use a cardinal
\(\kappa\geq2^{\aleph_0}\) with \(\kappa^{\aleph_0}\neq\kappa\), for
instance one of countable cofinality. Such an example, with
\(\kappa=\aleph_\omega\) and no hypothesis on the continuum, is
\cref{thm:p03-w03-p03-033-coding-theory}; Morley's footnote does not say
which example he had in mind. Neither theorem of this subsection uses total
transcendence, and neither is a substitute for
\cref{thm:p03-morley-saturation-existence}, which gives countably saturated
models in every uncountable power for totally transcendental theories,
including powers with \(\kappa^{\aleph_0}\neq\kappa\).
\end{hbwarning}

\subsection{Worked example and relationships}
\label{sec:p03-w03-p03-034-example}

\begin{hbexample}{def:p03-w03-p03-034-acf-example}{Algebraically closed fields under the results of this section}
Let \(\Sigma=\mathrm{ACF}_0\). By \cref{thm:p03-morley-field-example} it
is totally transcendental, with \(\alpha_T=2\), and categorical in every
uncountable power. For a countable subsystem \(A\), the canonical basis of
\(S(A)\) is countably infinite, matching
\cref{thm:p03-w03-p03-034-basis-power}: the formulas \(v_0=a\) for
\(a\in A\) and the polynomial equations over the field generated by \(A\)
are infinitely many inequivalent formulas. In every uncountable power
\(\kappa\) the unique model, the algebraically closed field of
transcendence degree \(\kappa\), is saturated by
\cref{thm:p03-morley-original-categoricity}, hence universal and homogeneous
for \(\mathcal N(T)\) by \cref{thm:p03-w03-p03-034-universal-homogeneous}:
an elementary map between two subsets of the field of power less than
\(\kappa\) extends to an automorphism of the whole field, and every
subsystem of a model of power at most \(\kappa\) embeds in it. The
countable algebraic closure of \(\mathbb Q\) is not saturated, since it
omits the transcendental type over \(\varnothing\); this agrees with the
failure of \(\aleph_0\)-categoricity.
\end{hbexample}

The results of this section depend on the following parts of
\cref{sec:p03-n-a-006-01}: the normalization and amalgamation theorem
(\cref{thm:p03-morley-normalization}), the realization lemma
(\cref{thm:p03-morley-types}), the countability characterization
(\cref{thm:p03-morley-countability}) and the main theorem
(\cref{thm:p03-morley-original-categoricity}). The remaining footnoted
material of the article, namely the single-point variant of the rank, the
automorphism test for unique types, the comparison of total transcendence
with density of isolated types, the endomorphism consequence and
question~(5), and the attributions to Ehrenfeucht, Vaught and Scott, is
treated in \cref{sec:p03-w03-p03-033-morley-supplement}.
